% Cross-Model Confidence in the MCL Framework
% Framework draft v0.1 -- 2026-09-13
% Compile with pdfLaTeX.
\PassOptionsToPackage{unicode}{hyperref}
\PassOptionsToPackage{hyphens}{url}
\documentclass[
]{article}
\usepackage{xcolor}
\usepackage{amsmath,amssymb}
\setcounter{secnumdepth}{-\maxdimen} % section numbers are written manually in the headings
\usepackage{iftex}

\usepackage[T1]{fontenc}
\usepackage[utf8]{inputenc}
\usepackage{textcomp}
\usepackage{lmodern}

\pdfoutput=1

\usepackage[hidelinks]{hyperref}
\hypersetup{
  pdfauthor   = {Li Zhang},
  pdftitle    = {Cross-Model Confidence in the MCL Framework: Calibrated Uncertainty for Selective Judgment},
  pdfsubject  = {Preprint},
  pdfkeywords = {multi-model collaboration; confidence calibration; selective prediction; uncertainty; LLM-as-judge; Meta-Coordination Layer}
}
\providecommand{\tightlist}{%
  \setlength{\itemsep}{0pt}\setlength{\parskip}{0pt}}
\setlength{\emergencystretch}{3em}
\usepackage[margin=1in]{geometry}
\usepackage{booktabs}
\usepackage{graphicx}
\usepackage{microtype}

\begin{document}
\section{Cross-Model Confidence in the MCL Framework: Calibrated Uncertainty for Selective Judgment}\label{title}

\begin{center}\rule{0.5\linewidth}{0.5pt}\end{center}

\subsection{Title \& Author Information}\label{title-author-information}

\begin{itemize}
\tightlist
\item
  \textbf{Title}: Cross-Model Confidence in the MCL Framework: Calibrated Uncertainty for Selective Judgment
\item
  \textbf{Author}: Li Zhang
\item
  \textbf{Affiliation}: TuringCorp
\item
  \textbf{Version}: framework draft v0.1 (2026-09-13)
\end{itemize}

\begin{center}\rule{0.5\linewidth}{0.5pt}\end{center}

\subsection{Abstract}\label{abstract}

A single model answers a question and, if asked, reports a confidence about its own
answer. That signal is introspective: nothing outside the model checks it, and a model
that is wrong can be wrong with high stated confidence. This paper studies an
alternative produced by the Meta-Coordination Layer (MCL) framework, in which several
heterogeneous models generate and cross-examine independently and an arbitration step
selects among the resulting candidates, emitting a confidence value with every
decision.

We evaluate the framework through its two deployed instantiations --- a collaborative
generator and a collaborative judge --- on three benchmark families: two public
judgment benchmarks (JudgeBench, 614 scored judgments; ContextualJudgeBench, the full
official 8-split set, 1,991 scored pairs / 3,988 response orders) and an expert-level,
multi-domain rubric benchmark (ProfBench, 38 scored tasks). We ask whether cross-model
confidence is calibrated against outcomes; how much accuracy a confidence threshold
buys at what coverage; where the signal degrades; whether the decision is invariant to
presentation order, unlike a single model used as a judge; and whether
preference-based judgment separates systems differently from rubric scoring.

On the public judgment benchmarks the collaborative judge's raw accuracy ties with
DeepSeek V4 Flash called directly (JudgeBench 92.5\% vs 92.7\%; CJB 67.1\%
vs 65.4\%), while its confidence is monotone in outcome accuracy across four bands
(JudgeBench 99.6\% / 94.0\% / 84.1\% / 67.7\%; CJB 83.3\% / 76.4\% / 63.6\% / 55.4\%).
Applying a threshold turns this into selective judgment: on JudgeBench, gating at
80\% confidence retains 76.1\% of judgments at 97.4\% accuracy, and at 90\% retains
46.1\% at 99.6\%. On the rubric benchmark, the collaborative system scores 63.3
against 57.4 for DeepSeek V4 Flash called directly and 55.6 for the official o3 draft on the
same 38 tasks under one judging pipeline; in a rubric-blind side-by-side comparison the
collaborative draft was selected in 36 of 38 tasks, and exchanging the two candidates
between the option positions kept the same pick in 10 of 10 cases.

The paper's central negative result is that the usefulness of the confidence signal is
conditional on how verifiable the underlying judgment is. On CJB splits where the
correct judgment is checkable, gating at 70\% confidence retains 93.6\% of decisions at
83.0\% accuracy; on splits whose two candidates are deliberately near-tied, the same
threshold retains 79.8\% at 62.6\% and the separation between correct and incorrect
decisions nearly vanishes. We report calibration, order-invariance and instrument
comparisons together with the comparisons in which the collaborative system does not
come out ahead.

\subsection{Keywords}\label{keywords}

multi-model collaboration; confidence calibration; selective prediction;
uncertainty quantification; LLM-as-judge; Meta-Coordination Layer

\subsection{1 Introduction}\label{introduction}

\subsubsection{1.1 Problem and background}\label{problem}

For expert questions --- a research question, a technical analysis, a business case ---
there is frequently no answer key. Evaluation therefore shifts from \emph{is this
correct} to \emph{which of these is better}, and any deployed system must additionally
answer a third question that no benchmark score covers: \emph{how much should this
judgment be trusted?}

A single model supplies a verdict and, when prompted, a confidence. The two come from
the same forward pass, so the confidence is a report about the model's own internal
state rather than evidence about the task. In practice this is visible as
high-confidence errors, and it is the reason a single-model confidence cannot be used
directly as an accept/reject gate.

The MCL framework takes a different route. Independent models produce candidates; the
candidates are cross-examined; an arbitration step selects among them and attaches a
confidence to the selection. The confidence is then a summary of a comparison between
concrete, externally produced candidates rather than an introspection, and the
comparison can be re-run with the candidates exchanged.

\subsubsection{1.2 Research questions and hypotheses}\label{questions}

\begin{table}[h]
\centering
\small
\begin{tabular}{@{}p{0.06\linewidth}p{0.52\linewidth}p{0.34\linewidth}@{}}
\toprule
& \textbf{Hypothesis} & \textbf{Status in this draft} \\
\midrule
H1 & Cross-model confidence is monotone in outcome accuracy. & Supported on two public judgment benchmarks and one rubric benchmark. \\
H2 & Confidence supports selective judgment with a quantified risk--coverage trade-off. & Supported; magnitude depends on benchmark. \\
H3 & The informativeness of the signal depends on how verifiable the judgment is. & Supported; the paper's central negative result. \\
H4 & A single model's self-reported confidence is not an equivalent substitute. & \emph{[pending]} --- planned measurement, no numbers reported yet. \\
H5 & The collaborative decision is invariant to presentation order, unlike a single-model judge. & Supported for the collaborative decision; a single-model judge shows a measurable order effect. \\
H6 & On questions without an answer key, preference-based judgment separates systems differently from rubric scoring. & Partially supported; the two instruments disagree and the disagreement is analysed rather than averaged. \\
\bottomrule
\end{tabular}
\end{table}

\subsubsection{1.3 Contributions and scope}\label{contributions}

\begin{enumerate}
\def\labelenumi{\arabic{enumi}.}
\tightlist
\item
  A calibration study of cross-model confidence on two public judgment benchmarks and
  one expert-level rubric benchmark, reported per confidence band with sample sizes.
\item
  A selective-judgment analysis that converts the calibration into an actionable
  trade-off: accuracy at a given threshold against the coverage it retains.
\item
  A boundary result: the signal is informative where the judgment is verifiable and
  degrades toward no information where the candidates are near-tied.
\item
  An order-invariance control for the collaborative decision, contrasted with the
  order sensitivity of a single model used as a judge.
\item
  An instrument comparison on questions without an answer key, in which rubric
  scoring and preference judgment are reported side by side and the reasons for their
  disagreement are analysed.
\end{enumerate}

The paper treats a collaborative generator and a collaborative judge as two sides of
one mechanism and studies them with the same reporting conventions.

\subsubsection{1.4 What this paper does not claim}\label{not-claim}

We do not claim that collaboration is more accurate than a single model in general: on
one public judgment benchmark the raw pick rate ties, and we report that as a result
rather than as a caveat. We do not claim that the confidence value replaces domain
review, or that a threshold transfers across task families --- the boundary result
above says the opposite. We do not claim that the rubric scores in this paper are
comparable to scores produced by a different judging pipeline; the calibration of our
pipeline against official labels is measured and reported so that the two scales are
kept apart.

\subsection{2 Related Work}\label{related-work}

\subsubsection{2.1 Multi-model collaboration and ensembling}\label{ensembles}

Ensembling and multi-agent debate improve accuracy on reasoning benchmarks, and the
mechanism is usually attributed to diversity among the members. Most of this literature
reports accuracy. This paper asks a complementary question about the same machinery:
what does the collaboration know about its own decisions, and can that knowledge be
turned into a gate.

\subsubsection{2.2 Confidence calibration}\label{calibration}

Calibration work generally measures whether a model's stated probability matches its
empirical accuracy, and treats overconfidence as a property of the model to be
corrected post hoc. We study a confidence that is produced by a comparison between
externally generated candidates rather than by introspection, and we measure it on the
same items for which accuracy is measured, so that the two axes can be read together.

\subsubsection{2.3 Selective prediction}\label{selective}

Selective prediction trades coverage for accuracy using a confidence score, and is
evaluated with risk--coverage curves. We adopt this framing because it states the
practical value of a confidence signal in the units a user cares about, and because it
makes the negative case explicit: when the curve is flat, the signal carries no
information for that task family.

\subsubsection{2.4 Models as judges}\label{llm-judge}

Using a model to judge outputs is now standard practice, and its known failure modes
include position bias and self-preference. Both matter directly here: the paper
measures the collaborative decision's sensitivity to position, and it compares
preference judgments produced by different instruments, including a single-model judge
that shares provenance with one of the candidates.

\subsection{3 The MCL Framework Under Study}\label{framework}

The Meta-Coordination Layer is a framework for coordinating heterogeneous models:
candidates are generated independently, cross-examined, and fused by a selection step,
with the whole procedure repeated at several quality--cost operating points. The
framework is instantiated in this paper through two deployed systems, described here at
the level of abstraction at which the framework itself is published.

\begin{itemize}
\tightlist
\item
  \textbf{Collaborative generator.} Given an open question, it returns a deliverable
  consisting of an answer body together with the reasoning the model states for it.
\item
  \textbf{Collaborative judge.} Given a task and two candidate answers, it returns
  which candidate is better, a confidence value between 0 and 100\%, and a short
  reason. The judge never sees a rubric or a score.
\end{itemize}

Both are exposed as OpenAI-compatible endpoints and were run against the benchmarks in
Section 4 through their public interfaces, using the same code paths that serve
production traffic. Internal configuration --- member selection, coordination depth and
parameterisation of the operating points --- is not part of this paper and is not
required to reproduce any number reported here, all of which are measurements of
externally observable behaviour.

\subsection{4 Experimental Setup and Reporting Conventions}\label{setup}

\subsubsection{4.1 Benchmarks}\label{benchmarks}

\begin{table}[h]
\centering
\small
\begin{tabular}{@{}llrl@{}}
\toprule
\textbf{Benchmark} & \textbf{What it measures} & \textbf{Size} & \textbf{Source} \\
\midrule
ProfBench & expert deliverables, scored criterion by criterion & 40 tasks, 4 domains & NVIDIA, public \\
JudgeBench & judging between two answers, objective items & 620 pairs & public \\
ContextualJudgeBench & judging between near-tied candidates, 8 splits & 2,000 pairs & Salesforce, public \\
\bottomrule
\end{tabular}
\end{table}

All three are third-party public benchmarks. We do not redistribute their contents; the
protocols are theirs and are followed verbatim.

\subsubsection{4.2 Protocols and aggregation}\label{protocols}

\begin{itemize}
\tightlist
\item
  \textbf{Rubric scoring.} Per-criterion Yes/No with the official prompt verbatim,
  temperature 0, top\_p 1, no output cap. Scores are \emph{task-level means}: each task
  is averaged first and the tasks are then averaged. Pooled per-criterion means differ
  from this convention by up to 1.7 points on our data, so the two are not
  interchangeable and only the task-level convention is reported.
\item
  \textbf{Judging.} First verdict per pair under the official protocol. For CJB,
  \emph{consistent accuracy} requires the same correct pick in both response orders;
  the random floor is 25\%.
\item
  \textbf{Exclusions.} Excluded items are disclosed rather than imputed: 2 of 40
  ProfBench tasks (no answer was produced), 6 of 620 JudgeBench pairs and 12 CJB orders
  (platform failures after retries).
\item
  \textbf{Judge calibration.} Our rubric pipeline is calibrated against the official
  per-criterion labels on the official reference draft: 74.0\% agreement, F1 0.761 on
  1,162 criteria. The residual per-domain bias is reported with the results, and
  absolute scores are not compared across pipelines.
\end{itemize}

\subsubsection{4.3 What is measured against what}\label{against}

Every comparison in this paper is between systems evaluated on the same items, under
the same protocol, with the same judging pipeline, and all reported margins are
accompanied by the counts they are computed from.

The single-model baseline throughout is \textbf{DeepSeek V4 Flash, called directly}
through its official API with no coordination: no second opinion, no cross-examination
and no arbitration. The confidence-elicitation run in Section~5.3 uses the same
endpoint and, since it is executed later, may be served by a later build of that
endpoint than the original runs; this is recorded as a limitation rather than
corrected.

\subsection{5 RQ1: Is Cross-Model Confidence Calibrated?}\label{rq1}

\subsubsection{5.1 Confidence bands against outcomes}\label{bands}

\begin{table}[h]
\centering
\small
\begin{tabular}{@{}lrrrr@{}}
\toprule
\textbf{Confidence band} & \multicolumn{2}{c}{\textbf{JudgeBench}} & \multicolumn{2}{c}{\textbf{CJB}} \\
\cmidrule(lr){2-3}\cmidrule(lr){4-5}
 & share & accuracy & share & accuracy \\
\midrule
$\geq$ 90\% & 45.6\% & 99.6\% & 19.8\% & 83.3\% \\
80--90\% & 29.7\% & 94.0\% & 37.3\% & 76.4\% \\
70--80\% & 13.2\% & 84.1\% & 29.7\% & 63.6\% \\
$<$ 70\% & 10.5\% & 67.7\% & 13.3\% & 55.4\% \\
\bottomrule
\end{tabular}
\end{table}

Monotonicity holds on both benchmarks. The two differ in where the mass sits, and that
difference is itself informative: on JudgeBench nearly half of the judgments are made
at $\geq$ 90\% confidence, while on CJB --- whose splits are designed to be near-tied
--- only a fifth are.

\subsubsection{5.2 Separation between correct and incorrect decisions}\label{separation}

Mean stated confidence is 85.8\% for correct and 65.2\% for incorrect judgments on
JudgeBench (a 20.6-point separation), and 80.0\% against 73.6\% on CJB (6.4 points).
The narrowing of the separation is the first sign of the boundary result developed in
Section 7.

\subsubsection{5.3 Self-reported confidence of a single model \emph{[pending]}}\label{single-model-confidence}

This subsection is reserved for the measurement of a single model's self-reported
confidence on the same judged items, elicited under the same protocol with a confidence
field added to its output. No numbers are reported here until that measurement is
complete. The comparison is pre-declared: the same items, the same protocol, the same
metrics, one model instead of the collaborative judge.

\subsection{6 RQ2: Selective Judgment}\label{rq2}

\subsubsection{6.1 Risk--coverage}\label{risk-coverage}

Gating a decision on its confidence converts calibration into a policy. Table~3 reports
the accuracy retained at three thresholds and the share of decisions each threshold
keeps.

\begin{table}[h]
\centering
\small
\begin{tabular}{@{}lrrr@{}}
\toprule
\textbf{Data} & \textbf{no gate} & \textbf{$\geq$ 70\%} & \textbf{$\geq$ 80\%} & \textbf{$\geq$ 90\%} \\
\midrule
JudgeBench & 92.5\% & 89.4\% / 95.4\% & 76.1\% / 97.4\% & 46.1\% / 99.6\% \\
CJB, checkable splits & 81.5\% & 93.6\% / 83.0\% & 77.6\% / 84.2\% & 34.7\% / 84.7\% \\
CJB, near-tied splits & 60.9\% & 79.8\% / 62.6\% & 36.5\% / 67.4\% & 4.9\% / 73.5\% \\
\bottomrule
\end{tabular}
\end{table}

Cells are \emph{coverage / accuracy}. The top row is the case the signal is for: a user
willing to set aside one judgment in ten buys roughly three points of accuracy. The
bottom row is the case it is not for: at the same threshold the gate adds under two
points, and pushing to 90\% leaves almost nothing to act on. Both are useful answers ---
the second says the two candidates really are close, which is a decision the user can
make on other grounds.

\subsubsection{6.2 Reading the gate as a product claim}\label{gate-product}

The operational statement of this section is deliberately narrow: a confidence value is
worth shipping when the curve it produces is not flat, and the curve must be measured
per task family rather than advertised once. \emph{[Section to be extended once the
single-model comparison of 5.3 is available, so that the two curves can be shown
together.]}

\subsection{7 RQ3: Where the Signal Holds and Where It Does Not}\label{rq3}

\subsubsection{7.1 Verifiability, not difficulty}\label{verifiability}

Splitting CJB by whether the correct judgment is checkable separates two regimes with
the same model and the same confidence scale. On checkable splits the gate is
efficient; on near-tied splits the difference between the mean confidence of correct
and incorrect decisions shrinks to 3.2 points and the gate buys almost nothing. The
relevant variable is therefore not how hard the item is but whether the comparison can
be settled by inspection at all.

\subsubsection{7.2 Order invariance}\label{order}

\begin{itemize}
\tightlist
\item
  The collaborative judge, re-run with the two candidates exchanged between the option
  positions, kept the same pick in 10 of 10 sampled tasks.
\item
  A single model used as a judge --- DeepSeek V4 Flash, over 76 head-to-head
  comparisons --- picked the candidate in the second position at a rate 7.9 points
  higher than the candidate in the first.
\end{itemize}

The contrast is the point: the collaborative decision is stable under an intervention
that measurably moves a single-model judge.

\subsection{8 RQ4: Preference and Rubric Are Different Instruments}\label{rq4}

\subsubsection{8.1 Two measurements of the same comparison}\label{two-measurements}

On ProfBench, scored criterion by criterion on the same 38 tasks by the same pipeline,
the collaborative generator scores 63.3 against 57.4 for DeepSeek V4 Flash called
directly and 55.6 for the official o3 reference draft. In a separate comparison the
same tasks were put to the collaborative judge as a rubric-blind preference question
--- the task and two drafts, no rubric, no scores --- and the collaborative draft was
chosen in 36 of 38 tasks, at a median confidence of 85.9\%.

\subsubsection{8.2 Instrument hierarchy and its limits}\label{instrument-limits}

The two instruments do not measure the same thing, and the paper reports the hierarchy
rather than choosing the flattering one:

\begin{itemize}
\tightlist
\item
  rubric scoring separates the systems by a few points;
\item
  a single-model judge sharing provenance with one candidate produces an approximately
  even split, and in the cases where it disagrees with the rubric, it favours its own
  provenance;
\item
  the collaborative judge, which never sees the rubric, picks one side almost
  unanimously, and does so without a position effect.
\end{itemize}

None of these is ground truth. The preference result is produced by a system of our
own, which is exactly the kind of self-preference risk documented for model judges; the
rubric result comes from a pipeline whose calibration is measured but not perfect. The
defensible reading is the one this section states: \emph{on questions without an answer
key the measured difference between systems depends on which instrument is asked, and
the instruments disagree in a structured way.}

\subsection{9 Discussion and Limitations}\label{limitations}

\begin{enumerate}
\def\labelenumi{\arabic{enumi}.}
\tightlist
\item
  \textbf{Accuracy parity is a result here, not a caveat.} On JudgeBench the
  collaborative judge's raw accuracy ties with DeepSeek V4 Flash called directly. The paper's
  claim is about the uncertainty axis, not about the pick rate.
\item
  \textbf{Every preference measurement in this paper was produced by one of the systems
  under comparison.} The order control addresses position, not provenance. An
  independent third-party judging instrument is required before the preference margin
  can be read as a property of the drafts rather than of the instrument.
\item
  \textbf{The rubric pipeline is calibrated, not identical, to the official one}
  (74.0\% agreement, F1 0.761, residual per-domain bias). Cross-pipeline score
  comparisons are therefore excluded, and only same-pipeline margins are reported.
\item
  \textbf{Sample sizes are small at the task level.} ProfBench contributes 38 scored
  tasks; the order-invariance control covers 10 of them. We report these as
  engineering-grade evidence and do not perform significance testing.
\item
  \textbf{The confidence scale is family-dependent.} Bands transfer across the two
  public judgment benchmarks in ordering but not in level, and the paper shows rather
  than assumes the difference.
\item
  \textbf{The single-model confidence comparison is pending.} Until it is measured, H4
  is stated as an open question and the abstract reports only measured values.
\item
  \textbf{Model versions drift.} Benchmarks were run at different times against hosted
  endpoints. The baseline is DeepSeek V4 Flash throughout; the confidence-elicitation
  run of Section~5.3 executes later than the original runs and may therefore be served
  by a later build of the same endpoint, and the raw accuracy of that run is reported
  alongside its confidence so that any drift is visible rather than assumed away.
\end{enumerate}

\subsection{10 Conclusion}\label{conclusion}

A collaborative system's most useful output may not be a better answer but a calibrated
statement about how much the answer should be trusted. On two public judgment
benchmarks the collaborative judge's accuracy ties with DeepSeek V4 Flash called directly while
its confidence is monotone in outcome accuracy and supports a quantified
risk--coverage trade-off; on an expert rubric benchmark the collaborative generator
leads by a few points under rubric scoring and is preferred almost unanimously under a
rubric-blind comparison. The same data show the limit of the claim: where the two
candidates are genuinely close, the confidence signal stops separating right from
wrong, and the honest reading of a low value is not \emph{the system failed} but
\emph{the choice belongs to criteria outside the answers}.

\subsection{Appendix A: Data Availability}\label{data-availability}

Aggregate results, the derived tables used in this paper and the scripts that produce
them are published at
\url{https://github.com/TuringCorp-net/Paper-Calibrated-Confidence-MCL}. The three
benchmarks are third-party and are referenced rather than redistributed:

\begin{itemize}
\tightlist
\item
  ProfBench --- \url{https://huggingface.co/datasets/nvidia/ProfBench}
\item
  JudgeBench --- as used in the published protocol
\item
  ContextualJudgeBench --- Salesforce, official 8-split release
\end{itemize}

Raw per-call logs, provider billing records and internal configuration are not
published.

\begin{thebibliography}{9}
\tightlist
\item
  NVIDIA. \emph{ProfBench: Multi-Domain Rubrics requiring Professional Knowledge to
  Answer and Judge.} arXiv:2510.18941, 2026.
\item
  Salesforce. \emph{ContextualJudgeBench.} 2026.
\item
  JudgeBench. \emph{Judging between candidate answers across knowledge, reasoning,
  mathematics and code items.} 2026.
\item
  Zhang, L. \emph{MCL conceptual framework.} Preprint, 2026.
  \url{https://doi.org/10.6084/m9.figshare.30645869}
\item
  Zhang, L. \emph{MCL-Team: minimal multi-model collaboration.} Preprint, 2026.
  \url{https://doi.org/10.6084/m9.figshare.32824352}
\end{thebibliography}

\end{document}
